\documentclass{article}
\usepackage{amsmath,amssymb,mathtools}
\usepackage[margin=1in]{geometry}
\allowdisplaybreaks
\begin{document}

\section*{Executable RNDDPC}
\begin{align}
\min_{\mathbf{u}(k),\mathbf{x}(k)}\quad &\sum_{i=0}^{N-1}\left(\|x(i+1\mid k)-r(i+1\mid k)\|_Q^2+\|u(i\mid k)\|_R^2\right) \tag{50a}\\
\text{s.t.}\quad
&\hat{\mathcal{R}}^z(i+1\mid k)=\mathcal{M}_{\mathrm{ABHJ}}\left(\hat{\mathcal{R}}^z(i\mid k)\times\mathcal{Z}_u\times\mathcal{Z}_{\epsilon}\times\mathcal{Z}_{\vartheta}\right)+\mathcal{Z}_{\sigma}, \tag{50b}\\
&\hat{\mathcal{R}}(i+1\mid k)=\mathcal{M}_C\hat{\mathcal{R}}^z(i+1\mid k)+\mathcal{Z}_{\varrho}, \tag{50c}\\
&\hat{\mathcal{H}}_l(i+1\mid k)\geq\mathcal{X}_l,\qquad \hat{\mathcal{H}}_u(i+1\mid k)\leq\mathcal{X}_u, \tag{50d}\\
&x(i+1\mid k)\geq\hat{\mathcal{H}}_l(i+1\mid k),\qquad x(i+1\mid k)\leq\hat{\mathcal{H}}_u(i+1\mid k), \tag{50e}\\
&u(i\mid k)\in\mathcal{U}, \tag{50f}\\
&x(0\mid k)=x(k), \tag{50g}\\
&\hat{\mathcal{R}}(0\mid k)=\langle x(k),0\rangle. \tag{50h}
\end{align}
The stage constraints (50b)--(50f) run over the prediction horizon, $i=0,\ldots,N-1$. The time $k$ is the current discrete time step. This is equation (50) of Li et al., \emph{Robust Nonlinear Data-Driven Predictive Control for Mixed Vehicle Platoons via Koopman Operator and Reachability Analysis}, arXiv:2502.16466v2, physical PDF page 17. Only the executable interval version is imposed; the earlier reachable-set membership constraints are not additional constraints of this formulation.

\subsection*{Decision variables, fixed inputs, and state constraints}
The decision sequences and their component domains are
\[
\begin{aligned}
\mathbf{u}(k)&=\{u(0\mid k),u(1\mid k),\ldots,u(N-1\mid k)\}, &u(i\mid k)&\in\mathbb{R},\\
\mathbf{x}(k)&=\{x(1\mid k),x(2\mid k),\ldots,x(N\mid k)\}, &x(i+1\mid k)&\in\mathbb{R}^{2n}.
\end{aligned}
\]
Here $N$ is the prediction-horizon length. The measured $x(k)$, the reference, weights, state and input limits, learned lifting function, offline data, and uncertainty descriptions are inputs to the online problem. The reachable sets and interval endpoints in (50) are computed by the displayed set propagation and interval construction, not independent state trajectories.

The definitions accompanying (46), on physical PDF pages 15--16, are
\[
\begin{gathered}
r(i+1\mid k)=\left[r_1^\top(k),r_2^\top(k),\ldots,r_n^\top(k)\right]^\top\in\mathbb{R}^{2n},\\
r_i(k)=\left[s_i^*(k),v_i^*(k)\right]^\top,\qquad i\in\{1,\ldots,n\},\\
Q=\operatorname{diag}\left(Q_x,\xi Q_x,\ldots,\xi^{(n-1)}Q_x\right)\in\mathbb{R}^{2n\times2n},\qquad R\in\mathbb{R},\\
Q_x=\operatorname{diag}(\rho_s,\rho_v),\qquad 0<\xi\leq1,\\
\tilde{x}_{\max}=[\tilde{s}_{\max},\tilde{v}_{\max}]^\top,\\
\mathcal{X}=\left\{x(k)\in\mathbb{R}^{2n}\ \middle|\ |x(k)-r(k)|\leq\mathbf{1}_n\otimes\tilde{x}_{\max}\right\},\\
\mathcal{U}=\left\{u(k)\in\mathbb{R}\ \middle|\ |u(k)|\leq u_{\max}\right\}.
\end{gathered}
\]
The quantities $s_i^*(k)$ and $v_i^*(k)$ are desired spacing and velocity. The weights $\rho_s$ and $\rho_v$ penalize spacing and velocity deviations; $\xi$ decreases the penalties for HDVs farther from the CAV. The symbols $\tilde{s}_{\max}$ and $\tilde{v}_{\max}$ are deviation limits, and $u_{\max}$ is the maximum allowable CAV input. The vector $\mathbf{1}_n$ is the $n$-dimensional column vector of ones, and $\otimes$ is the Kronecker product. The vectors $\mathcal{X}_l$ and $\mathcal{X}_u$ in (50d) are the lower and upper state bounds of $\mathcal{X}$. The paper does not give additional general sign restrictions on $\rho_s,\rho_v,R$ in this definition; none are added here.

\subsection*{Physical variables and bounded adverse factors}
There is one leading CAV, indexed by $1$, followed by $n-1$ HDVs; the head vehicle has index $0$. The source definitions on physical PDF pages 5--6 are
\[
\Omega=\{1,2,\ldots,n\},\qquad \Omega_C=\{1\},\qquad \Omega_H=\{2,\ldots,n\},
\]
\[
s_i(k)=p_{i-1}(k)-p_i(k),\qquad x_i(k)=[s_i(k),v_i(k)]^\top,\qquad i\in\Omega,
\]
\begin{equation}
x(k)=\left[x_1^\top(k),x_2^\top(k),\ldots,x_n^\top(k)\right]^\top\in\mathbb{R}^{2n}. \tag{11}
\end{equation}
Here $p_i$ and $v_i$ are position and velocity. The scalar input is CAV acceleration; the disturbance is head-vehicle velocity, and the attack is an additive non--time-delay input attack:
\[
u(k)=u_1(k)\in\mathbb{R},\qquad \epsilon(k)=v_0(k)\in\mathbb{R},\qquad \vartheta(k)=\vartheta_1(k)\in\mathbb{R},\qquad \omega(k)\in\mathbb{R}^{2n}.
\]
The underlying nonlinear dynamics are left abstract by the data-driven formulation. The bounded-signal premises printed in (13) are
\begin{equation}
|\epsilon(k)|\leq\epsilon_{\max},\qquad |\vartheta(k)|\leq\vartheta_{\max},\qquad |\omega(k)|\leq\omega_{\max}, \tag{13}
\end{equation}
where $\epsilon_{\max},\vartheta_{\max}\in\mathbb{R}$ and $\omega_{\max}\in\mathbb{R}^{2n}$. Noise is unknown and is not directly measurable during data collection.

\subsection*{Interval construction and set operations}
The interval construction replacing the earlier set constraints is
\begin{equation}
\hat{\mathcal{H}}(i+1\mid k)=\operatorname{interval}\bigl(\hat{\mathcal{R}}(i+1\mid k)\bigr)
=[\hat{\mathcal{H}}_l(i+1\mid k),\hat{\mathcal{H}}_u(i+1\mid k)]. \tag{49}
\end{equation}
Its endpoints are obtained by (5), below, applied to the propagated zonotope. The preliminary definitions on physical PDF page 4 use $n$ as a generic set dimension, rather than specifically the number of platoon vehicles.

An interval set is defined by
\[
\mathcal{H}=\left\{x_{\mathcal{H}}\in\mathbb{R}^n\ \middle|\ \underline{x}_{\mathcal{H}i}\leq x_{\mathcal{H}i}\leq\overline{x}_{\mathcal{H}i},\quad\forall i=1,\ldots,n\right\}
=[\underline{\mathcal{H}},\overline{\mathcal{H}}],
\]
\[
\underline{\mathcal{H}}=[\underline{x}_{\mathcal{H}1},\underline{x}_{\mathcal{H}2},\ldots,\underline{x}_{\mathcal{H}n}],\qquad
\overline{\mathcal{H}}=[\overline{x}_{\mathcal{H}1},\overline{x}_{\mathcal{H}2},\ldots,\overline{x}_{\mathcal{H}n}].
\]
For $c_{\mathcal{Z}}\in\mathbb{R}^n$, $\gamma_{\mathcal{Z}}\in\mathbb{N}$, and
\[
G_{\mathcal{Z}}=[g_{\mathcal{Z}}^{(1)},g_{\mathcal{Z}}^{(2)},\ldots,g_{\mathcal{Z}}^{(\gamma_{\mathcal{Z}})}]\in\mathbb{R}^{n\times\gamma_{\mathcal{Z}}},
\]
the zonotope definition is
\begin{equation}
\mathcal{Z}=\langle c_{\mathcal{Z}},G_{\mathcal{Z}}\rangle
=\left\{x\in\mathbb{R}^n\ \middle|\ x=c_{\mathcal{Z}}+\sum_{i=1}^{\gamma_{\mathcal{Z}}}\beta^{(i)}g_{\mathcal{Z}}^{(i)},\ -1\leq\beta^{(i)}\leq1\right\}. \tag{3}
\end{equation}
For compatible dimensions the printed operations are
\[
L\mathcal{Z}=\langle Lc_{\mathcal{Z}},LG_{\mathcal{Z}}\rangle,\qquad L\in\mathbb{R}^{m\times n},
\]
\[
\mathcal{Z}_1+\mathcal{Z}_2=\left\langle c_{\mathcal{Z}_1}+c_{\mathcal{Z}_2},[G_{\mathcal{Z}_1},G_{\mathcal{Z}_2}]\right\rangle,
\qquad \mathcal{Z}_a=\langle c_{\mathcal{Z}_a},G_{\mathcal{Z}_a}\rangle\ (a=1,2),
\]
\begin{equation}
\mathcal{Z}_1\times\mathcal{Z}_2=
\left\langle\begin{bmatrix}c_{\mathcal{Z}_1}\\c_{\mathcal{Z}_2}\end{bmatrix},
\begin{bmatrix}G_{\mathcal{Z}_1}&0\\0&G_{\mathcal{Z}_2}\end{bmatrix}\right\rangle. \tag{4}
\end{equation}
The interval over-approximation is
\begin{equation}
\mathcal{H}=\operatorname{interval}(\mathcal{Z})=[c_{\mathcal{Z}}-\Delta g_{\mathcal{Z}},c_{\mathcal{Z}}+\Delta g_{\mathcal{Z}}],\qquad
\Delta g_{\mathcal{Z}}=\sum_{i=1}^{\gamma_{\mathcal{Z}}}|g_{\mathcal{Z}}^{(i)}|, \tag{5}
\end{equation}
where the absolute value is taken element-wise.

For $C_{\mathcal{M}}\in\mathbb{R}^{n\times m}$, $\gamma_{\mathcal{M}}\in\mathbb{N}$, and
\[
G_{\mathcal{M}}=[G_{\mathcal{M}}^{(1)},G_{\mathcal{M}}^{(2)},\ldots,G_{\mathcal{M}}^{(\gamma_{\mathcal{M}})}]\in\mathbb{R}^{n\times m\gamma_{\mathcal{M}}},
\]
a matrix zonotope is
\begin{equation}
\mathcal{M}=\langle C_{\mathcal{M}},G_{\mathcal{M}}\rangle
=\left\{X\in\mathbb{R}^{n\times m}\ \middle|\ X=C_{\mathcal{M}}+\sum_{i=1}^{\gamma_{\mathcal{M}}}\beta^{(i)}G_{\mathcal{M}}^{(i)},\ -1\leq\beta^{(i)}\leq1\right\}. \tag{6}
\end{equation}
The $X$ in this generic definition is a matrix element, not the offline state-data matrix defined below.

\subsection*{Lifting and the uncertainty-bearing predictor}
The state lifting and its domain, from physical PDF pages 8 and 11, are
\begin{equation}
\Psi(\chi(k))=\begin{bmatrix}z(k)\\u(k)\\\epsilon(k)\\\vartheta(k)\end{bmatrix},\qquad
z(k)=[\varphi_1(x(k)),\varphi_2(x(k)),\ldots,\varphi_{n_z}(x(k))]^\top\in\mathbb{R}^{n_z},\qquad n_z>2n. \tag{19}
\end{equation}
For the selected DNN lifting,
\begin{equation}
z(k)=\begin{bmatrix}x(k)\\g(k)\end{bmatrix}. \tag{26}
\end{equation}
The first $2n$ basis functions retain the physical state, and the remaining basis functions are learned:
\[
x(k)=[\varphi_1(x(k)),\ldots,\varphi_{2n}(x(k))]^\top,\qquad
g(k)=[\varphi_{2n+1}(x(k)),\ldots,\varphi_{n_z}(x(k))]^\top.
\]
The encoder defining $g(k)$ has $n_e$ fully connected layers,
\begin{equation}
o_e^{(p)}=\sigma_e^{(p)}\left(W_e^{(p)}o_e^{(p-1)}+b_e^{(p)}\right),\qquad p\in\{1,2,\ldots,n_e\}. \tag{27}
\end{equation}
Here $W_e^{(p)}\in\mathbb{R}^{n_p\times n_{p-1}}$, $b_e^{(p)}\in\mathbb{R}^{n_p}$, and $n_p$ is the layer width. The input and output are $o_e^{(0)}=x(k)$ and $o_e^{(n_e)}=g(k)$. ReLU activations are used for $p\in\{1,\ldots,n_e-1\}$; the final layer has no activation. These learned quantities are fixed offline, not online decision variables.

The matrices have domains
\[
[A\ B\ H\ J]\in\mathbb{R}^{n_z\times(n_z+3)},\quad
A\in\mathbb{R}^{n_z\times n_z},\quad B,H,J\in\mathbb{R}^{n_z},\quad C\in\mathbb{R}^{2n\times n_z}.
\]
They represent the system, control, disturbance, attack, and projection matrices, respectively. The model used to define the error and model sets is the equivalent Koopman-based model on physical PDF page 13:
\begin{equation}
\begin{cases}
z(k+1)=Az(k)+Bu(k)+H\epsilon(k)+J\vartheta(k)+\sigma(k),\\
x(k+1)=Cz(k+1)+\varrho(k).
\end{cases} \tag{36}
\end{equation}
This is the uncertainty-bearing model underlying the reachable-set construction, not an additional nominal equality imposed on the decision trajectory in (50). The modeling errors combine modeling inaccuracies and noise. Their assumed bounds are
\begin{equation}
|\sigma(k)|\leq\sigma_{\max},\qquad |\varrho(k)|\leq\varrho_{\max},\qquad
\sigma_{\max}\in\mathbb{R}^{n_z},\quad\varrho_{\max}\in\mathbb{R}^{2n}. \tag{37}
\end{equation}
The error zonotopes are
\begin{equation}
\begin{aligned}
\sigma(k)\in\mathcal{Z}_{\sigma}
&=\langle c_{\mathcal{Z}_{\sigma}},G_{\mathcal{Z}_{\sigma}}\rangle
=\langle0_{n_z\times1},\operatorname{diag}(\sigma_{\max})\rangle,\\
\varrho(k)\in\mathcal{Z}_{\varrho}
&=\langle c_{\mathcal{Z}_{\varrho}},G_{\mathcal{Z}_{\varrho}}\rangle
=\langle0_{2n\times1},\operatorname{diag}(\varrho_{\max})\rangle.
\end{aligned} \tag{38}
\end{equation}
The bounds are obtained from statistical characteristics of prediction errors evaluated on data. Appendix E gives the residual expressions used for this purpose, with the learned matrices and lifting known:
\begin{equation}
\begin{cases}
\sigma(k)=z(k+1)-Az(k)-Bu(k)-H\epsilon(k)-J\vartheta(k),\\
\varrho(k)=x(k+1)-Cz(k+1).
\end{cases} \tag{E.1}
\end{equation}
No simulation-specific values or probability-coverage percentages replace the bounded-error premise (37).

\subsection*{Offline data defining the model sets}
The collection phase applies persistently exciting control, disturbance, and attack sequences of length $T+1$ under noisy conditions. The input, disturbance, and attack signals are manually specified and the state is measured. The source data definitions are
\begin{align}
U&=[u(1),u(2),\ldots,u(T+1)]\in\mathbb{R}^{1\times(T+1)}, \tag{23a}\\
E&=[\epsilon(1),\epsilon(2),\ldots,\epsilon(T+1)]\in\mathbb{R}^{1\times(T+1)}, \tag{23b}\\
F&=[\vartheta(1),\vartheta(2),\ldots,\vartheta(T+1)]\in\mathbb{R}^{1\times(T+1)}, \tag{23c}\\
X&=[x(1),x(2),\ldots,x(T+1)]\in\mathbb{R}^{2n\times(T+1)}, \tag{23d}\\
Z&=[z(1),z(2),\ldots,z(T+1)]\in\mathbb{R}^{n_z\times(T+1)}. \tag{24}
\end{align}
The matrices used in the construction are
\begin{align}
U_-&=[u(1),u(2),\ldots,u(T)]\in\mathbb{R}^{1\times T}, \tag{25a}\\
E_-&=[\epsilon(1),\epsilon(2),\ldots,\epsilon(T)]\in\mathbb{R}^{1\times T}, \tag{25b}\\
F_-&=[\vartheta(1),\vartheta(2),\ldots,\vartheta(T)]\in\mathbb{R}^{1\times T}, \tag{25c}\\
X_-&=[x(1),x(2),\ldots,x(T)]\in\mathbb{R}^{2n\times T}, \tag{25d}\\
X_+&=[x(2),x(3),\ldots,x(T+1)]\in\mathbb{R}^{2n\times T}, \tag{25e}\\
Z_-&=[z(1),z(2),\ldots,z(T)]\in\mathbb{R}^{n_z\times T}, \tag{25f}\\
Z_+&=[z(2),z(3),\ldots,z(T+1)]\in\mathbb{R}^{n_z\times T}. \tag{25g}
\end{align}
The distinct symbols $U$ and $\mathcal{U}$ denote a data matrix and an admissible-input set; $X$ and $\mathcal{X}$ likewise denote a data matrix and a state-constraint set.

\subsection*{Matrix-zonotope model construction and its premises}
Lemmas 1 and 2 require both
\[
\begin{bmatrix}Z_-\\U_-\\E_-\\F_-\end{bmatrix}
=\left[Z_-^\top\ U_-^\top\ E_-^\top\ F_-^\top\right]^\top
\quad\text{and}\quad [Z_+]
\quad\text{to have full row rank}.
\]
Under these premises, the printed model sets are
\begin{align}
\mathcal{M}_{\mathrm{ABHJ}}&=(Z_+-\mathcal{M}_{\sigma})\begin{bmatrix}Z_-\\U_-\\E_-\\F_-\end{bmatrix}^{\dagger}, \tag{39}\\
\mathcal{M}_C&=(X_+-\mathcal{M}_{\varrho})[Z_+]^{\dagger}, \tag{40}\\
\mathcal{M}_{\sigma}&=\left\langle C_{\mathcal{M}_{\sigma}},[G_{\mathcal{M}_{\sigma}}^{(1)},G_{\mathcal{M}_{\sigma}}^{(2)},\ldots,G_{\mathcal{M}_{\sigma}}^{(n_zT)}]\right\rangle, \tag{41}\\
\mathcal{M}_{\varrho}&=\left\langle C_{\mathcal{M}_{\varrho}},[G_{\mathcal{M}_{\varrho}}^{(1)},G_{\mathcal{M}_{\varrho}}^{(2)},\ldots,G_{\mathcal{M}_{\varrho}}^{(2nT)}]\right\rangle. \tag{42}
\end{align}
The dagger denotes the Moore--Penrose pseudoinverse. The complete generator construction is
\begin{align}
C_{\mathcal{M}_{\sigma}}&=[c_{\mathcal{Z}_{\sigma}}\ \cdots\ c_{\mathcal{Z}_{\sigma}}]\in\mathbb{R}^{n_z\times T}, \tag{43a}\\
G_{\mathcal{M}_{\sigma}}^{(1+(i-1)T)}&=[g_{\mathcal{Z}_{\sigma}}^{(i)}\quad 0_{n_z\times(T-1)}]\in\mathbb{R}^{n_z\times T}, \tag{43b}\\
G_{\mathcal{M}_{\sigma}}^{(k+(i-1)T)}&=[0_{n_z\times(k-1)}\quad g_{\mathcal{Z}_{\sigma}}^{(i)}\quad0_{n_z\times(T-k)}]\in\mathbb{R}^{n_z\times T}, \tag{43c}\\
G_{\mathcal{M}_{\sigma}}^{(T+(i-1)T)}&=[0_{n_z\times(T-1)}\quad g_{\mathcal{Z}_{\sigma}}^{(i)}]\in\mathbb{R}^{n_z\times T}, \tag{43d}\\
C_{\mathcal{M}_{\varrho}}&=[c_{\mathcal{Z}_{\varrho}}\ \cdots\ c_{\mathcal{Z}_{\varrho}}]\in\mathbb{R}^{2n\times T}, \tag{44a}\\
G_{\mathcal{M}_{\varrho}}^{(1+(j-1)T)}&=[g_{\mathcal{Z}_{\varrho}}^{(j)}\quad0_{2n\times(T-1)}]\in\mathbb{R}^{2n\times T}, \tag{44b}\\
G_{\mathcal{M}_{\varrho}}^{(k+(j-1)T)}&=[0_{2n\times(k-1)}\quad g_{\mathcal{Z}_{\varrho}}^{(j)}\quad0_{2n\times(T-k)}]\in\mathbb{R}^{2n\times T}, \tag{44c}\\
G_{\mathcal{M}_{\varrho}}^{(T+(j-1)T)}&=[0_{2n\times(T-1)}\quad g_{\mathcal{Z}_{\varrho}}^{(j)}]\in\mathbb{R}^{2n\times T}. \tag{44d}
\end{align}
Here $g_{\mathcal{Z}_{\sigma}}^{(i)}$ and $g_{\mathcal{Z}_{\varrho}}^{(j)}$ are the respective columns of the generator matrices in (38). The source quantifier ranges are
\[
i\in\{1,2,\ldots,n_z\},\qquad j\in\{1,2,\ldots,2n\},\qquad k\in\{2,3,\ldots,T-1\}.
\]
The $k$ in this generator construction is a column-position index, not the current online time.

Appendix A identifies the error sequences underlying these sets as
\begin{align}
O_-&=[\sigma(1),\sigma(2),\ldots,\sigma(T)]\in\mathbb{R}^{n_z\times T}, \tag{A.1}\\
\Gamma_-&=[\varrho(1),\varrho(2),\ldots,\varrho(T)]\in\mathbb{R}^{2n\times T}, \tag{A.2}\\
Z_+&=[A\ B\ H\ J]\begin{bmatrix}Z_-\\U_-\\E_-\\F_-\end{bmatrix}+O_-, \tag{A.3}\\
X_+&=CZ_++\Gamma_-. \tag{A.4}
\end{align}
These error sequences are not measurable. The matrix zonotopes over-approximate the system matrices consistent with the data and bounded errors; Appendix B uses
\[
[A\ B\ H\ J]\in\mathcal{M}_{\mathrm{ABHJ}},\qquad [C]\in\mathcal{M}_C.
\]

\subsection*{Reachability inputs and scope of the selected replacement}
The recursion (45) is exactly the pair of propagation rows retained as (50b)--(50c). Its accompanying definitions, on physical PDF page 15, are
\[
\mathcal{Z}_u=\langle u(i\mid k),0\rangle,\qquad
\mathcal{Z}_{\epsilon}=\langle\epsilon(k),\epsilon_{\max}\rangle,\qquad
\mathcal{Z}_{\vartheta}=\langle0,\vartheta_{\max}\rangle.
\]
The set $\hat{\mathcal{R}}^z(i+1\mid k)$ is the over-approximated reachable set of the lifted state $z(i+1\mid k)$, and $\hat{\mathcal{R}}(i+1\mid k)$ is that of the physical state $x(i+1\mid k)$. Appendix B assumes that the exact and over-approximated lifted reachable sets originate from the same initial set. The input zonotope depends on the decision $u(i\mid k)$; the disturbance zonotope is centered at the measured head-vehicle velocity, and the attack zonotope is centered at zero. The zero-attack and zero-noise assumptions of the earlier nominal controllers are not imposed on (50).

\paragraph{Printed scope and omissions.}
The printed initialization (50h) specifies only $\hat{\mathcal{R}}(0\mid k)$, whereas (50b) needs $\hat{\mathcal{R}}^z(0\mid k)$. Neither the selected block nor the linked initialization passages give a formula for this initial lifted set. It is not supplied here. Also, the source calls $\mathcal{Z}_{\epsilon}=\langle\epsilon(k),\epsilon_{\max}\rangle$ a set derived from (13), although (13) bounds the absolute disturbance about zero, not its variation about $\epsilon(k)$. Both printed expressions are retained with their original scopes; no centered-deviation bound is substituted for (13). Finally, (50e) is interval membership and does not impose membership in the original zonotope. The superseded zonotope-membership row is not reinstated.

\paragraph{Scope of the appendix guarantees.}
Appendices C and D state recursive feasibility and closed-loop constraint-satisfaction results for problem (48), not a separate theorem for (50). Assumption 1 in Appendix C requires $\mathcal{X}$ to be one-step robustly control invariant under some admissible backup input $u_b\in\mathcal{U}$. The closed-loop result additionally presumes disturbances, attacks, and modeling errors bounded by $\mathcal{Z}_{\epsilon},\mathcal{Z}_{\vartheta},\mathcal{Z}_{\sigma},\mathcal{Z}_{\varrho}$ and feasibility at the stated time steps. These are premises of the cited guarantees, not extra terminal constraints printed in (50).

\subsection*{Online use}
At each time step the problem is solved with the current measured state. The resulting sequences are
\[
\mathbf{u}^*(k)=[u^*(0\mid k),u^*(1\mid k),\ldots,u^*(N-1\mid k)],\qquad
\mathbf{x}^*(k)=[x^*(1\mid k),x^*(2\mid k),\ldots,x^*(N\mid k)].
\]
Only the first input $u^*(0\mid k)$ is applied to the CAV; the time is incremented and the measured state is updated. Algorithm 1 lists pre-collected $(U,E,F,X)$, constraints $(\mathcal{X},\mathcal{U})$, weights $(Q,R)$, bounded sets $(\mathcal{Z}_{\epsilon},\mathcal{Z}_{\vartheta},\mathcal{Z}_{\omega})$, horizon $N$, and total steps $N_s$ as inputs, and runs while $0\leq k\leq N_s$. It lists $\mathcal{Z}_{\omega}$ without a general defining formula; (50) instead uses the modeling-error zonotopes (38). No additional noise-zonotope construction is inserted into the selected optimization problem.

\end{document}
